% ============================================================================
% §3.3.2 "Nozzle exit profiles" — REVISED v2 (2026-07-22, after user review).
% Drop-in replacement for 04_cases_and_diagnostics.tex lines 170–322.
% Scope tightened to what De/256 supports: sensitivity study, not mechanism
% separation; no local-gradient/instability claims; velocity-based measures
% qualified; C_M = convective momentum-flux coefficient; discrete audit now
% quoted ON the comparison grid De/256 (convention validated against the
% logged D/8 and D/512 hole audits). Change log: section_3_3_2_review_notes.md
% ============================================================================
\subsection{Nozzle exit profiles}
\label{subsec:uej-inlet-profiles}

A finite-thickness velocity profile is prescribed at the nozzle exit as the
primary inlet condition. Within the nozzle orifice, the axial velocity
is defined by the shifted hyperbolic tangent function
\begin{equation}
    f(r)=\frac{1}{2}
    \left[1-\tanh\left(\frac{r-r_0}{\delta}\right)\right],
    \qquad r_0=R_e-3\delta,
    \label{eq:uej-tanh-profile}
\end{equation}
where \(R_e=D_e/2\) and \(\delta=0.01D_e=\SI{254}{\micro\meter}\). The
velocity is prescribed as \(u_x(r)=u_e f(r)\), with zero transverse
components. Locating the centre of the transition at \(R_e-3\delta\) reduces
the axial velocity to approximately \(0.0025u_e\) at the orifice edge.

The static pressure is uniform across the orifice, and the temperature and
density follow from a constant total temperature:
\[
    p(r)=p_e,
    \qquad
    T(r)=T_0-\frac{u_x^2(r)}{2c_p},
    \qquad
    \rho(r)=\frac{p_e}{R\,T(r)},
\]
so the radial variations in temperature and density are determined by
\(f(r)\), while the static pressure remains uniform. The resulting effective
discharge coefficient,
\(C_{d,\mathrm{eff}}=\dot m/(\rho_e u_e\pi R_e^2)\simeq0.879\), is a
property of the prescribed profile, not an experimentally measured discharge
coefficient. The profile is applied only for \(r<R_e\); the remainder of the
upstream boundary is specified in \S\ref{subsec:uej-computational-setup}.

The sensitivity of the computed near field to the prescribed nozzle-exit
velocity profile is assessed using the baseline profile and three additional
cases. The comparison examines how the mean shock-cell structure varies with
the prescribed profile shape, nominal shear-layer thickness, and integral
mass and convective momentum fluxes. Since several profile properties vary
simultaneously, the cases are treated as a sensitivity study rather than as a
strict separation of the underlying mechanisms.

The first additional case uses a top-hat profile, \(f_{\mathrm{TH}}=1\) for
\(r<R_e\) and zero otherwise, representing the idealised limit of an abrupt
velocity change at the nozzle lip; its continuous analytic
\(C_{d,\mathrm{eff}}\) is unity. The second case belongs to the shifted-tanh
family of Eq.~\eqref{eq:uej-tanh-profile} with
\(\delta=\SI{100}{\micro\meter}\) (\(\delta/D_e\simeq0.004\)). The third
case places the maximum velocity gradient at the nozzle lip:
\begin{equation}
    f_{\mathrm{WT}}(r)
    =\tanh\left(\frac{R_e-r}{a}\right),
    \qquad a=\SI{200}{\micro\meter}.
    \label{eq:uej-wall-tanh-profile}
\end{equation}

The prescribed profiles are characterised by their vorticity thickness,
half-velocity location, \SI{99}{\percent} thickness, velocity-profile
integral measures, and flux coefficients. Here,
\(\Delta U=u_e-u_\infty=u_e\) for the nominally quiescent ambient, while
\(y_{50}\) and \(\delta_{99}\) are measured inward from the nozzle lip.
The profile measures used here are defined by
\begin{subequations}
\label{eq:uej-profile-measures}
\begin{gather}
    \delta_\omega
    =\frac{\Delta U}{\max\lvert\partial u_x/\partial r\rvert},
    \\[0.25em]
    \begin{aligned}
    f(R_e-y_{50})&=\frac{1}{2},
    &\qquad
    f(R_e-\delta_{99})&=0.99,
    \end{aligned}
    \\[0.25em]
    \begin{aligned}
    \delta^*
    &=\int_0^{R_e}\left[1-f(r)\right]\,\mathrm{d}r,
    &\qquad
    \theta_0
    &=\int_0^{R_e}f(r)\left[1-f(r)\right]\,\mathrm{d}r,
    \end{aligned}
    \\[0.25em]
    H=\frac{\delta^*}{\theta_0},
    \\[0.25em]
    C_M=\frac{\displaystyle\int_{A_e}\rho u_x^2\,\mathrm{d}A}
             {\rho_e u_e^2\pi R_e^2}.
\end{gather}
\end{subequations}
The quantities \(\delta^*\), \(\theta_0\), and \(H\) are
velocity-profile measures evaluated using locally planar definitions. They
include neither density weighting nor the cylindrical area factor \(2\pi r\)
and are used only to characterise the imposed profile shapes. The coefficient
\(C_M\) is the normalised axial convective momentum-flux coefficient. It
excludes the exit-pressure and viscous contributions to thrust. Since the
same exit pressure is imposed in all four cases, it provides a consistent
measure of their profile-dependent convective momentum flux. Analytically,
\(\theta_0\simeq\delta/2\) for the shifted-tanh family and
\(\theta_0=a(1-\ln 2)\simeq0.31a\) for the wall-tanh profile.

The thin shifted-tanh and wall-tanh profiles are matched in their nominal
analytic vorticity thickness,
\(\delta_\omega=\SI{200}{\micro\meter}\), and have nearly identical
\(\delta_{99}\simeq\SI{530}{\micro\meter}\). Their gradient locations,
integral thicknesses, shape factors, and flux coefficients nevertheless
differ. Similar mean shock-cell measures for the pair would be consistent
with sensitivity to their common nominal thickness scale, whereas trends
ordered by \(C_{d,\mathrm{eff}}\) or \(C_M\) across the four profiles would
be consistent with sensitivity to the integral flux deficit. Neither
comparison identifies a unique causal mechanism because the complete profile
shapes remain different. The tabulated quantities are analytic properties of
the imposed profiles; the evolved shear layer is not assumed to retain them.

\begin{table}[H]
\centering
\caption{Velocity-profile measures and integrated flux coefficients of the
four nozzle exit profiles used in the sensitivity study. Values are analytic
properties of the imposed profiles; \(C_M\) is the normalised axial
convective momentum-flux coefficient.}
\label{tab:uej-exit-profile-integrals}
\small
\setlength{\tabcolsep}{4.2pt}
\renewcommand{\arraystretch}{1.16}
\begin{tabular}{@{}ccccccccc@{}}
\toprule
\multirow[c]{2}{*}{Profile} &
\multirow[c]{2}{*}{\shortstack{Scale\\(\si{\micro\meter})}} &
\multicolumn{4}{c}{Thickness measures (\si{\micro\meter})} &
\multirow[c]{2}{*}{\(H\)} &
\multicolumn{2}{c}{Integrated flux coefficients} \\
\cmidrule(lr){3-6}\cmidrule(lr){8-9}
& & \(\delta_\omega\) & \(y_{50}\) & \(\delta^*\) & \(\theta_0\) &
& \(C_{d,\mathrm{eff}}\) & \(C_M\) \\
\midrule
Top hat & -- & -- & -- & 0 & 0 & -- & 1.0000 & 1.0000 \\
Wall tanh & \(a=200\) & 200 & 110 & 139 & 61 & 2.26 & 0.9755 & 0.9669 \\
Shifted tanh (thin) & \(\delta=100\) & 200 & 300 & 300 & 50 & 6.02 & 0.9513 & 0.9445 \\
Shifted tanh (baseline) & \(\delta=254\) & 508 & 762 & 762 & 127 & 6.02 & 0.8789 & 0.8623 \\
\bottomrule
\end{tabular}
\end{table}

Figure~\ref{fig:uej-profiles-overview} compares the
four profiles in Table~\ref{tab:uej-exit-profile-integrals} with two additional
profile shapes used in the separate shear-layer comparisons of
\S\ref{sec:uej-profile-sensitivity}.

\begin{figure}[H]
    \centering
    \includegraphics[width=\textwidth]{profiles_overview.pdf}
    \caption[Nozzle exit velocity profiles]{Nozzle exit velocity profiles used
    for the inlet sensitivity calculations. The first two panels show all six
    profiles over the full radius and near the nozzle lip; the third shows the
    finite radial gradients of the four smooth profiles with nonsingular lip
    gradients. The Pohlhausen--Blasius and one-seventh-power profiles are used
    only in the separate profile-shape comparisons of
    \S\ref{sec:uej-profile-sensitivity}.}
    \label{fig:uej-profiles-overview}
\end{figure}

The four-profile calculations use
\(\Delta x=D_e/256=\SI{99.2}{\micro\meter}\). At this spacing, the nominal
\(\delta_\omega=\SI{200}{\micro\meter}\) spans only about two cells and
\(\delta_{99}\) about five cells. Cell-centred quadrature on the comparison
grid reproduces the orifice area and the analytic
\(C_{d,\mathrm{eff}}\) and \(C_M\) values of all four profiles to within
\SI{0.008}{\percent}. For the three finite-thickness profiles, both flux
coefficients agree with their analytic values to within
\SI{0.001}{\percent}; on a \(D_e/512\) grid, the corresponding errors do not
exceed \SI{0.003}{\percent}. Although the integrated fluxes are accurately
represented, the local velocity gradients and their instability development
are not sufficiently resolved for quantitative assessment. Conclusions are
therefore restricted to the sensitivity of the mean shock-cell system to the
discretely imposed exit profiles.
